% problem-000332 4 氮化锂层状晶体
\section*{4. 氮化锂层状晶体}
\textit{下列为分问。}

\subsection*{4-1}
分别画出 $\mathrm { L i _ { 2 } N ^ { - } }$ −层的结构，最⼩重复单位以及 $\mathrm { . } \mathrm { a - L i _ { 3 } N }$ 的晶胞。

\subsection*{4-2}
确定该晶体的结构基元、点阵型式以及 $\mathrm { N ^ { 3 - } }$ −的Li+离⼦配位数。

\subsection*{4-3}
计算 $\mathrm { L i _ { 2 } N }$ −层Li—N中间距离。

\subsection*{4-4}
通过计算说明α-Li_{3}N晶体导电的原因。假如 $\mathrm { N ^ { 3 - } }$ −作六⽅最密堆积(hcp)，指出Li+离⼦占据空隙类型及占据百分数，回答该结构的 $\mathrm { L i _ { 3 } N }$ 能否导电，简述理由。

\subsection*{4-5}
⾦属Li在⾼温下与 $\mathrm { T i O _ { 2 } }$ 反应可⽣成多种晶型的 $\mathrm { L i T i _ { 2 } O _ { 4 } }$ ，其中尖晶⽯型 $\mathrm { L i T i _ { 2 } O _ { 4 } }$ 为⾯⼼⽴⽅结构。 $\mathrm { O ^ { 2 \mathrm { - } } }$ −作⽴⽅最密堆积 $\left( \mathrm { c c p } \right)$ ，若Li+离⼦有序地占据四⾯体空隙，Ti4+和 $\mathrm { T i ^ { 3 + } }$ 离⼦占据⼋⾯体空隙。已知晶胞参数 $\mathbf { \partial } _ { \cdot } a \mathbf { \ } =$ 840.5 pm，晶体密度 $D = 3 . 7 3 0 \mathrm { g c m ^ { - 3 } }$ 。请分别计算晶胞中所含原⼦数⽬和 $\mathrm { O ^ { 2 - } }$ 离⼦半径的最⼤可能值。若不计其他离⼦，给出以Li+离⼦为顶点的最⼩单位。
